\documentclass[10pt,a4paper]{amsart}
\usepackage[margin=19mm]{geometry}
\usepackage{amsmath,amssymb,mathtools}
\usepackage[scr=rsfs,cal=euler]{mathalfa}
\usepackage{xfrac}
\usepackage[unicode,hidelinks,bookmarks=false]{hyperref}
\newcommand{\ie}{i.e.\ }
\newcommand{\cf}{cf.\ }
\newcommand{\resp}{respectively\ }
\newcommand{\Subsection}{\subsection}
\providecommand{\nobreakdash}{\nobreak\hbox{-}}
\setlength{\emergencystretch}{2em}
\multlinegap0pt
\usepackage{amssymb}
\newtheorem{lemma}{Lemma}
\newcommand{\BB}{\mathbf{B}}
\newcommand{\PPP}{\mathbb{P}}
\newcommand{\TT}{\mathbf{T}}
\DeclareMathOperator{\inter}{int}
\newcommand{\p}{\mathbf{p}}
\newcommand{\q}{\mathbf{q}}
\title{Density bounds for unit ball packings relative to their outer parallel domains}
\author{K\'aroly Bezdek, Zsolt L\'angi}
\date{Source: arXiv:2401.00645v2 (CC BY 4.0)}
\begin{document}
\maketitle

\noindent\emph{Source and licence.} Density bounds for unit ball packings relative to their outer parallel domains. Version arXiv:2401.00645v2 (CC BY 4.0), distributed by arXiv under the Creative Commons Attribution 4.0 International licence, http://creativecommons.org/licenses/by/4.0/, as recorded in the arXiv licence metadata for this exact version and shown on the abstract page https://arxiv.org/abs/2401.00645v2. Full licence notice: \texttt{licenses/arxiv-2401.00645-CC-BY-4.0.txt}.\par\smallskip

\noindent\textit{Editorial scope.} Counted unit: the article's lemma lem:constcurv1 (source0.tex lines 394--396). The article's complete original proof of it is delivered with it (source0.tex lines 398--404, one \texttt{proof} environment of 7 lines). No mathematics is written, repaired or re-derived: the statement and the proof are the article's own lines, in the article's own notation, and no retained line is substituted or re-flowed.  No further source-local context is needed: every cross-reference of every retained line resolves inside the two delivered spans.  Nothing was omitted from any retained span, so the package carries no omission record.  No retained line contains a citation, so the wrapper carries no bibliography and no bibliography entry.  No retained line contains a figure environment, an image-inclusion command or drawing code, so no image is delivered and the package declares no asset dependency.  Editorial formatting only, in addition: an AMS \texttt{amsart} preamble that restates the article's own declarations and macros used by the retained lines, namely the environments \texttt{lemma} on separate counters, together with the standard packages those lines need.  The retained environments print in this document as Lemma 1 in order of appearance, whereas the article prints the same items with the numbering of its own full document.  The complete native source of the article as distributed in the arXiv e-print for this exact version is bundled unchanged as \texttt{sources/152310/source0.tex}.
\begin{lemma}\label{lem:constcurv1}
Let $L$ be a line through $\q$, and let $R_p$ be a half line starting at a point $\p$ of $L \setminus \inter (\BB)$ such that $R_p$ is perpendicular to $L$. For any $s >0 $ (on the sphere for any $0 < s < \frac{\pi}{2})$, let $\p(s)$ denote the point of $R_p$ whose distance from $\p$ is $s$. Furthermore, for any $0 < s_1 < s_2$ (on the sphere for any $0 < s_1 < s_2 < \frac{\pi}{2})$, let $\TT(s_1,s_2)$ denote the triangle with vertices $\q, \p(s_1)$ and $\p(s_2)$. Then $\rho(\BB,\TT(s_1,s_2))$ and $\hat{\rho}(\BB,\TT(s_1,s_2))$ are non-increasing functions of $s_1$ and $s_2$.
\end{lemma}

\begin{proof}
Let $0  < s_1 < s_2 < s_2'$ be given (on the sphere assume also that $s_2' < \frac{\pi}{2}$). Without loss of generality, assume that the ratio of the angles of $\TT(s_1,s_2)$ and $\TT(s_1,s_2')$ at $\q$ is rational. Then we can dissect $\TT(s_1,s_2')$ into triangles $\TT_1, \TT_2, \ldots, \TT_k$, all having $\q$ as a vertex such that they have equal angles at $\q$, denoted by $\gamma$, any two consecutive members share a side, and $\p(s_2)$ is a vertex of some triangles $\TT_i$ and $\TT_{i+1}$. Let $\varphi : \PPP \to \PPP$ be the rotation around $\q$ with angle $\gamma$ that moves $\p(s_1)$ closer to $\p(s_2)$. Then, for any value of $i$, $\varphi(T_i) \subseteq T_{i+1}$, $\varphi(\TT_i \cap \BB) = \TT_{i+1} \cap \BB$ and $\varphi(\TT_i \cap \BB_{\lambda}) = \TT_{i+1} \cap \BB_{\lambda}$.
This proves that $\rho(\BB,\TT(s_1,s_2)) \geq \rho(\BB,\TT(s_1,s_2'))$, implying that $\rho(\BB,\TT(s_1,s_2))$ is a non-increasing function of $s_2$. To prove that it is a non-increasing function of $s_1$, we may apply a similar argument.

To verify the monotonicity property of $\hat{\rho}(\BB,\TT(s_1,s_2))$, note that there is a unique value $\bar{s}$ such that $\p(s) \in \BB_{\lambda}$ if and only if $s \leq \bar{s}$. Furthermore, if $0 < s_1 < s_2 \leq \bar{s}$, then $\hat{\rho}  (\BB,\TT(s_1,s_2)) =1$. Thus, to prove that $\hat{\rho}(\BB,\TT(s_1,s_2))$ is a non-increasing function of $s_2$, it is sufficient to consider values $\bar{s} \leq s_1 < s_2 < s_2'$. But in this case, repeating the argument in the previous paragraph, we obtain
$\hat{\rho}(\BB,\TT(s_1,s_2)) \geq \hat{\rho}(\BB,\TT(s_1,s_2'))$, showing that the function is non-increasing in the second variable. To prove the monotonicity in the first variable, we can repeat the same argument.
\end{proof}

\end{document}
